\documentclass[twocolumn]{iopjournal}

\usepackage{amsmath,amssymb,amsfonts}

\fancyhfoffset[L]{0mm}
\setlength{\marginparwidth}{0pt}
\setlength{\marginparsep}{0pt}

\renewcommand{\articletype}[1]{%
  {\vspace*{-8mm}\noindent \Large \sf Preprint}\par
  \vspace*{6mm}\noindent{\scriptsize \sf{\bfseries \MakeUppercase{#1}}}}

\fancyhead[L]{}
\fancyhead[R]{}

\hypersetup{pdftitle={The SST Covariant Completion in Its Quasi-Static Limit}, pdfauthor={Stephen Nelson}, pdfsubject={Gravity Field Dynamics (GFD)}, pdfkeywords={scalar-tensor transition, covariant action, AQUAL, constitutive law, radial acceleration relation, Tully-Fisher, lensing, proof}}

\begin{document}

\articletype{Supplementary Material}

\title{The SST Covariant Completion in Its Quasi-Static Limit}

\author{Stephen Nelson}

\affil{Field Dynamics Research}

\ead{stephen@fielddynamics.org}

\begin{abstract}
The scalar-tensor transition (SST) is the four-term covariant action
of Gravity Field Dynamics, with kinetic function
$F(y) = y - 2\sqrt y + 2\ln(1+\sqrt y)$. We derive, exactly, what it
gives in the quasi-static weak-field limit in which galaxies are
observed. We read the action as the framework does: $\Phi$ is the
Newtonian potential of the metric, the curvature term carries the
tensor sector, and $\Phi$'s response is carried by $F$. Lemma~1 states
that reading as an exact identity for the static Lagrangian. From it
the curvature term's equation gives $\Psi = \Phi$, so $\gamma = 1$ at
leading order, and the $\Phi$-equation is
$\nabla\cdot[\mu(|\nabla\Phi|/a_0)\nabla\Phi] = 4\pi G\rho$ with
$\mu(x) = x/(x+1)$, equivalently the capacity form with
$f(x) = (21x+13)/(x+1)$. Minimal coupling of matter gives the source.
In spherical symmetry the field equation integrates exactly to
$g_s = g_t^2/(a_0 + g_t)$, whose forward solver and mass decoder are
exact inverses. We derive the Newtonian and deep limits, the baryonic
Tully--Fisher relation, the bound on the TetB share, the transition
ring, the phantom density and well-posedness, and, from $\gamma = 1$,
light deflection $4GM/(bc^2)$ in the Newtonian field and
$2\pi\sqrt{GMa_0}/c^2$ in the deep field, with lensing mass equal to
dynamical mass.
\end{abstract}

\keywords{scalar-tensor transition, covariant action, AQUAL, constitutive law, radial acceleration relation, Tully--Fisher, lensing, proof}

% ----------------------------------------------------------------------

\section*{Setting}

The SST is the action~\cite{NelsonL3,NelsonL3supp,NelsonII}
\begin{equation}
\begin{aligned}
S = \int d^4x\,\sqrt{-g}\,\Big[&\frac{R}{16\pi G}
 - \frac{a_0^2}{8\pi G}\,F(y)\\
&- \tfrac14 e^{-2\Phi/c^2}F_{\mu\nu}F^{\mu\nu}
 + \mathcal L_{\mathrm{matter}}\Big],
\end{aligned}
\label{eq:sst}
\end{equation}
with $y = |\nabla\Phi|^2/a_0^2$ and
$F(y) = y - 2\sqrt y + 2\ln(1 + \sqrt y)$. $F$ is the unique energy of
the constitutive law $\mu(x) = x/(x+1)$, $F'(y) = \mu(\sqrt y)$,
$F(0) = 0$~\cite{NelsonPaqual}; $\mu$ is derived from the exchange
symmetry of the two roots of the tetrad equation~\cite{NelsonL2,NelsonPexch}.
We use the framework's reading of the action, as stated in the
archive:
\begin{itemize}
\item[(i)] $\Phi$ is the Newtonian potential of the metric, ``not a new
auxiliary field; the partition modifies its response, not its
identity''~\cite{NelsonL3}, with the torsion share $1 - \mu$ carried
inside $F$~\cite{NelsonL3supp}.
\item[(ii)] The curvature term carries the tensor sector, so
gravitational waves travel at $c$, and $\Phi$'s response is carried by
$F$, the scalar sector ``decoupling from the metric
dynamics''~\cite{NelsonL3supp}. Lemma~1 states this exactly.
\item[(iii)] Matter is minimally coupled to the metric, and
$e^{-2\Phi/c^2} \to 1$ in the static weak field~\cite{NelsonRAR}.
\item[(iv)] $\mu$ applies to the net field of a galaxy, not to each
mass; a bound system much smaller than its host, such as the solar
system, is pure TetA, that is, general
relativity~\cite{NelsonColes,NelsonTetA}.
\end{itemize}
\textit{Naming.} ``SST'' names the action; the spherical, static
relation $g_s = g_t^2/(a_0 + g_t)$ is its spherical field
equation~\cite{NelsonColes}. The documentation calls the capacity
function $f(x) = (21x+13)/(x+1)$ ``the field equation''~\cite{NelsonDocs};
Theorem~2 shows that the two are one equation.

\paragraph{Contributions and prior work.} The AQUAL field equation, its
spherical reduction, the Newtonian and deep limits, the Tully--Fisher
relation, well-posedness and the deep-field deflection are standard
results for the simple interpolating
function~\cite{BekensteinMilgrom1984,Milgrom1983,FamaeyMcGaugh2012}.
What this proof adds is their derivation from the SST's own terms, with
the framework's reading stated as an exact identity (Lemma~1) and
$\gamma = 1$ derived from it, together with the capacity form, the
exact solver--decoder pair, the TetB-share bound and the transition
ring, each checked by machine.

% ----------------------------------------------------------------------

\section*{The Field Equation}

Write the static weak-field metric as $g_{00} = -(1 + 2\Phi/c^2)$,
$g_{ij} = (1 - 2\Psi/c^2)\delta_{ij}$.

\paragraph{Lemma 1 (the reading, exactly).} \textit{The static part of
$\sqrt{-g}R/16\pi G$ at second order is, up to a divergence,
$(8\pi G)^{-1}[(\nabla\Psi)^2 - 2\nabla\Phi\cdot\nabla\Psi]$, which on
$\Psi = \Phi$ is the Newtonian energy $-(\nabla\Phi)^2/8\pi G =
-(a_0^2/8\pi G)\,y$. Reading (ii) is the statement that the static
Lagrangian is}
\begin{equation}
\begin{aligned}
\mathcal L_{\mathrm{st}} = \;&\frac{(\nabla\Psi)^2 - 2\nabla\Phi\cdot\nabla\Psi}{8\pi G}\\
&- \frac{a_0^2}{8\pi G}\big[F(y) - y\big] - \rho\Phi ,
\end{aligned}
\label{eq:lst}
\end{equation}
\textit{the curvature term supplying the Newtonian part $y$ of $F$. On
$\Psi = \Phi$ this is $-(a_0^2/8\pi G)F(y) - \rho\Phi$. The source term
is minimal coupling: a mass $m$ contributes
$-mc^2\sqrt{-g_{00}} = -mc^2 - m\Phi + O(\Phi^2/c^2)$.}

\paragraph{Proof.} The second-order expansion of $\sqrt{-g}R$ for this
metric is computed directly; its Euler--Lagrange equations are twice
those of $(\nabla\Psi)^2 - 2\nabla\Phi\cdot\nabla\Psi$. The rest is
substitution. \hfill $\square$

\paragraph{Theorem 1.} \textit{Varying Eq.~(\ref{eq:lst}) with respect
to $\Psi$ gives $\nabla^2(\Phi - \Psi) = 0$, so $\Psi = \Phi$ for
fields that vanish at infinity: $\gamma = 1$ at leading order. Varying
with respect to $\Phi$ and setting $\Psi = \Phi$ gives}
\begin{equation}
  \nabla\cdot\Big[\mu\Big(\frac{|\nabla\Phi|}{a_0}\Big)\nabla\Phi\Big]
  = 4\pi G\rho .
\label{eq:field}
\end{equation}

\paragraph{Proof.} The $\Psi$-equation is
$(8\pi G)^{-1}(-2\nabla^2\Psi + 2\nabla^2\Phi) = 0$. In the
$\Phi$-equation, the curvature term contributes $\nabla^2\Psi/4\pi G$
and the term $+(a_0^2/8\pi G)y$ contributes $-\nabla^2\Phi/4\pi G$;
these cancel on $\Psi = \Phi$. For the rest, $\partial F/\partial(\nabla\Phi) =
F'(y)\,2\nabla\Phi/a_0^2 = 2\mu\nabla\Phi/a_0^2$, and
$(a_0^2/8\pi G)(2/a_0^2) = 1/(4\pi G)$. The Euler--Lagrange equation is
$(4\pi G)^{-1}\nabla\cdot(\mu\nabla\Phi) - \rho = 0$, the Bekenstein--Milgrom form~\cite{BekensteinMilgrom1984}. \hfill $\square$

\paragraph{Theorem 2 (the capacity form).} \textit{With
$\sigma(x) = (x-1)/(x+1)$, the capacity function
$f(x) = 17 + 4\sigma(x) = (21x+13)/(x+1)$ takes the values $13, 17, 21$
at $x = 0, 1, \infty$, and $\mu = (f - 13)/8$. So
Eq.~(\ref{eq:field}) is
$\nabla\cdot[\tfrac18(f - 13)\nabla\Phi] = 4\pi G\rho$, the form of
Paper~VIII and the documentation.}

\paragraph{Proof.} $17(x+1) + 4(x-1) = 21x + 13$, and
$(21x + 13 - 13x - 13)/8 = x$. \hfill $\square$

% ----------------------------------------------------------------------

\section*{The Spherical Field Equation}

\paragraph{Theorem 3.} \textit{In spherical symmetry, with
$g_t = |\Phi'(r)|$ and $g_s = GM({<}r)/r^2$, Eq.~(\ref{eq:field})
integrates exactly to the conserved flux $r^2\mu(g_t/a_0)\,g_t = GM({<}r)$,
that is,}
\begin{equation}
  g_s = \frac{g_t^2}{a_0 + g_t},\qquad
  g_t = \tfrac12\Big[g_s + \sqrt{g_s^2 + 4g_s a_0}\,\Big],
\label{eq:spherical}
\end{equation}
\textit{the second being the unique positive root. With $g_t = v^2/r$,
the mass decoder}
\begin{equation}
  M_{\mathrm{dec}}(r) = \frac{r v^4}{G(a_0 r + v^2)}
\end{equation}
\textit{is the exact inverse of the forward solver.}

\paragraph{Proof.} For $\Phi = \Phi(r)$,
$\nabla\cdot(\mu\nabla\Phi) = r^{-2}(r^2\mu\Phi')'$, and
$4\pi r^2\rho = M'(r)$, so the divergence theorem gives the flux with no
curl term. Substituting $\mu = x/(x+1)$, $x = g_t/a_0$, gives the
quadratic $g_t^2 - g_s g_t - g_s a_0 = 0$, whose other root is
negative. The decoder is $r^2 g_s/G$ at $g_t = v^2/r$. \hfill $\square$

\medskip
\noindent
This is the relation fitted to the SPARC rotation curves with no
parameter in $\mu$~\cite{NelsonRAR}, solved forward in
Paper~X~\cite{NelsonX} and inverted in the decoding-mass
letter~\cite{NelsonDEC}.

% ----------------------------------------------------------------------

\section*{Limits and Scales}

\paragraph{Theorem 4.} \textit{(i) As $g_s \to \infty$,
$g_t = g_s + a_0 - a_0^2/g_s + O(g_s^{-2})$: Newtonian. (ii) As
$g_s \to 0$, $g_t = \sqrt{g_s a_0} + g_s/2 + O(g_s^{3/2})$. (iii) For a
bounded, isolated mass $M$, $v^4 \to GMa_0$ exactly as $r \to \infty$:
the baryonic Tully--Fisher relation (an external field modifies this,
as in AQUAL). (iv) The TetB share of the field is
$(1-\mu)g_t = a_0\mu(g_t/a_0) < a_0$. (v) The ring $x = 1$, where
$\mu = \tfrac12$, lies where $g_s = a_0/2$, at
$r_1 = \sqrt{2GM/a_0}$; at $r = \sqrt{GM/a_0}$, where $g_s = a_0$, the
field is at $x = \varphi = (1+\sqrt5)/2$ with $\mu = 1/\varphi$.
(vi) Outside a point mass the phantom density
$\rho_{\mathrm{ph}} = (4\pi G r^2)^{-1}\,d(r^2 g_t)/dr$ has leading term
$\sqrt{GMa_0}/(4\pi G r^2)$.}

\paragraph{Proof.} (i), (ii): expansions of the forward solver. (iii):
$v^4 = g_t^2 r^2$ with $g_s = GM/r^2$, and the limit of the solver.
(iv): $g_t/(1+x) = a_0 x/(1+x)$. (v): $x = 1$ gives
$g_s = 1/(1+1)\cdot a_0$; $g_s = a_0$ gives $x^2 = x + 1$. (vi):
$r^2 g_s = GM$ is constant outside the mass, so the residual of the
Newtonian flux is $r^2 g_t$, which grows as $\sqrt{GMa_0}\,r$.
\hfill $\square$

\paragraph{Theorem 5 (well-posedness).} \textit{The energy density
$F(|p|^2/a_0^2)$ is strictly convex in $p = \nabla\Phi$ for $p \neq 0$,
with Hessian eigenvalues $2\mu$ across and $2(x\mu)'$ along the field,
whose ratio $(x\mu)'/\mu = 1 + (1-\mu)$ lies in $[1, 2]$. So
Eq.~(\ref{eq:field}) is elliptic away from $\nabla\Phi = 0$, and its
solution is unique for given boundary data or for fields decaying at
infinity. Existence for AQUAL is due to Bekenstein and
Milgrom~\cite{BekensteinMilgrom1984}.}

\paragraph{Proof.} The energy-functional proof~\cite{NelsonPaqual};
the eigenvalues are $2x/(x+1)$ and $2x(x+2)/(x+1)^2$. \hfill $\square$

% ----------------------------------------------------------------------

\section*{The Single Potential}

\paragraph{Theorem 6.} \textit{By Theorem~1 the weak-field metric
carries one potential in both components,
$g_{00} = -(1 + 2\Phi/c^2)$, $g_{ij} = (1 - 2\Phi/c^2)\delta_{ij}$, so the
lensing potential is $\int(\Phi + \Psi)\,dz = 2\int\Phi\,dz$: light is
deflected by the same $\Phi$ that sets the dynamics. In the Born
approximation the deflection at impact parameter $b$ is
$\theta = (2/c^2)\int g_\perp\,dz$. A Newtonian point mass gives
$4GM/(bc^2)$; the deep field $g_t = \sqrt{GMa_0}/r$ gives the constant
$2\pi\sqrt{GMa_0}/c^2$; and for a spherical lens the lensing mass equals
the dynamical mass $M_{\mathrm{dyn}}(r) = g_t r^2/G$. In a galaxy the
factor $e^{-2\Phi/c^2}$ differs from $1$ by about
$2|\Phi|/c^2 \sim 10^{-6}$, so the fourth term does not enter the
static dynamics.}

\paragraph{Proof.} The deflection is $((1+\gamma)/c^2)\int
\nabla_\perp\Phi\,dz$ with $\nabla_\perp\Phi = g_t\,b/r$; the integrals
$\int b\,r^{-3}dz = 2/b$ and $\int b\,r^{-2}dz = \pi$ give the two
values. A general-relativistic lens with mass profile $M_{\mathrm{dyn}}$
has $g = GM_{\mathrm{dyn}}/r^2 = g_t$, the same integrand.
\hfill $\square$

\medskip
\noindent
This is Letter~3's statement that $\gamma = 1$ ``because the single
potential $\Phi$ governs both metric components'', and Paper~I's
single-standard postulate~\cite{NelsonI}, here derived from the
curvature term's $\Psi$-equation. The deep-field deflection is
Milgrom's~\cite{Milgrom1983}.

\paragraph{Remark (the reading in (ii)).} The static weak-field part of
$R/16\pi G$ is $(8\pi G)^{-1}[(\nabla\Psi)^2 - 2\nabla\Phi\cdot\nabla\Psi]$,
which on $\Psi = \Phi$ is the Newtonian energy $-(\nabla\Phi)^2/8\pi G$.
If $F$ were added to it in full, rather than $F - y$ as in Lemma~1, the
equations would give $\Psi = \Phi$ and $\nabla\cdot[(1+\mu)\nabla\Phi]
= 4\pi G\rho$. The SST's field equation, Eq.~(\ref{eq:field}), is the
one in which $\Phi$'s response is carried by $F$ alone, as the archive
states; Lemma~1 is that statement made exact.

% ----------------------------------------------------------------------

\section*{Scope}

Theorems~1--6 are exact consequences of the SST, read as in Lemma~1, at
leading order in the quasi-static weak field, and Theorems~3--4 are
what the SPARC rotation curves, the Tully--Fisher relation and the mass
decoder test. The limit uses the static frame, in which $|\nabla\Phi|^2$
is defined; as a spacetime scalar, $y$ requires $\Phi$ to be a scalar
field, and its covariant form for time-dependent and strong fields is
given in the companion proof (\texttt{proofs/sst-covariant-time-leg/}):
$y = c^4a_\mu a^\mu/a_0^2$, with $a$ the acceleration of the tetrad's
time leg, and $\Phi$ the lapse. The archive records what else remains open~\cite{NelsonColes}:
the metric field equations beyond leading order, including the
stress--energy of the $F$ sector; $\mathcal L_{\mathrm{matter}}$ beyond
its weak-field form; the documents describe the fourth term
differently (disformal and needed for lensing; conformal with a
cosmological role; an auxiliary vector), and the dependence of the
fine-structure constant on the potential that it implies is not
discussed; and the match from the locked bimetric theory to $\mu$ is
open. At the radius $\sqrt{GM/a_0}$ quoted as the processing boundary in
Paper~I, the Newtonian acceleration equals $a_0$, and there
$x = \varphi$ and $\mu = 1/\varphi$; the ring $\mu = \tfrac12$ is at
$\sqrt{2GM/a_0}$ (Theorem~4(v)).

% ----------------------------------------------------------------------

\begin{thebibliography}{99}

\bibitem{NelsonL3}
S.~Nelson,
\textit{Gravity Field Dynamics: Field Closure of the Tetrad Frame},
Letter~3 (2026).

\bibitem{NelsonL3supp}
S.~Nelson,
\textit{Covariant completion}, supplementary note to Letter~3 (2026).

\bibitem{NelsonII}
S.~Nelson,
\textit{A Covariant Theory Without Dark Matter}, Paper~II (2026),
doi:10.5281/zenodo.18360354.

\bibitem{NelsonPaqual}
S.~Nelson,
\textit{The energy functional of the constitutive law} (2026),
\texttt{proofs/aqual-energy-functional/}.

\bibitem{NelsonL2}
S.~Nelson,
\textit{The Other Half of Gravity: Torsion and the Two-Chamber
Structure of the Tetrad Frame}, Letter~2 (2026).

\bibitem{NelsonPexch}
S.~Nelson,
\textit{The constitutive law from exchange symmetry} (2026),
\texttt{proofs/}\linebreak\texttt{constitutive-law-exchange-symmetry/}.

\bibitem{NelsonRAR}
S.~Nelson,
\textit{The Radial Acceleration Relation from a Zero-Parameter
Covariant Action} (2026).

\bibitem{NelsonColes}
S.~Nelson,
\textit{Gravity Field Dynamics --- Coles Notes} (2026), \S\S3, 5.4, 12.

\bibitem{NelsonTetA}
S.~Nelson,
\textit{Einstein's equation as the TetA partition} (2026).

\bibitem{NelsonDocs}
S.~Nelson,
\textit{Documentation: The Field Equation; Covariant Completion}
(2026), \texttt{docs/02-gravitation/}.

\bibitem{NelsonX}
S.~Nelson,
\textit{One Equation, Two Scales: Tully--Fisher Evolution and the
Hubble Constant}, Paper~X (2026).

\bibitem{NelsonDEC}
S.~Nelson,
\textit{Decoding Mass from Observed Rotation Curves} (2026),
doi:10.5281/zenodo.18721752.

\bibitem{NelsonI}
S.~Nelson,
\textit{The Gravitational Constitutive Law}, Paper~I (2026),
doi:10.5281/zenodo.18332595.

\bibitem{BekensteinMilgrom1984}
J.~Bekenstein and M.~Milgrom,
``Does the missing mass problem signal the breakdown of Newtonian
gravity?'', \textit{Astrophys. J.} \textbf{286} (1984) 7.

\bibitem{FamaeyMcGaugh2012}
B.~Famaey and S.~S.~McGaugh,
``Modified Newtonian Dynamics (MOND): observational phenomenology and
relativistic extensions'', \textit{Living Rev. Relativ.} \textbf{15}
(2012) 10.

\bibitem{Milgrom1983}
M.~Milgrom,
``A modification of the Newtonian dynamics as a possible alternative
to the hidden mass hypothesis'', \textit{Astrophys. J.} \textbf{270}
(1983) 365.

\end{thebibliography}

\end{document}
